\begin{partbacktext}
\part{The theory of \texorpdfstring{$\mathcal{I}$}{I}-good singularities}
\noindent
This part develops the comparison theory of singularities that forms the technical core of the book. The basic objects are positive-mass $\theta$-psh potentials, their $P$- and $\mathcal I$-singularity types, and the envelopes introduced in \cref{chap:envelopes}. The central theme is that the ordinary pluripotential-theoretic equivalence relation and the equivalence relation detected by multiplier ideals are different in general, but coincide on a large and geometrically meaningful class of singularities: the $\mathcal I$-good ones.

The first chapter of the part, \cref{chap:comp}, introduces the general language for comparing singularities. It studies the $P$- and $\mathcal I$-preorders and defines the pseudometric $d_S$ on positive-mass potentials. This pseudometric is not restricted to $\mathcal I$-good potentials; rather, it is the ambient convergence tool used to compare singularity types. The chapter proves the main continuity and convergence properties of $d_S$, including its compatibility with envelopes and with changes of the background form.

In \cref{chap:Igood} we introduce $\mathcal I$-good singularities. The main structural result is the characterization theorem \cref{thm:charIgoodasclosure}, which relates $\mathcal I$-goodness to approximation by analytic singularities, equality between non-pluripolar mass and volume, and the coincidence of the $P$- and $\mathcal I$-envelopes. The chapter also develops the mixed volume formalism needed later and proves the asymptotic Riemann--Roch formula \cref{thm:DXmain1}, which identifies the non-pluripolar volume of an $\mathcal I$-good current with the asymptotic growth of multiplier-ideal sections.

\cref{chap:trace} develops the trace operator. Its purpose is to give a meaningful restriction of a quasi-psh singularity to a possibly singular subvariety, in a way compatible with multiplier ideals and birational modifications. The trace operator is accompanied by analytic Bertini-type results and by restricted-volume formulae. These results allow the singularity theory developed in the previous chapters to be used inductively on subvarieties.

\cref{chap:testcurve} studies test curves, namely concave decreasing families of model potentials. It proves the Ross--Witt Nystr\"om correspondence between test curves and geodesic rays and develops the operations on test curves needed later. The chapter also introduces the $\mathcal I$-model version of the theory, which will be used in the non-Archimedean applications.

\cref{chap:Okounkov} develops partial Okounkov bodies for singular metrics. These convex bodies extend the classical Newton--Okounkov construction to the setting of psh metrics and record the valuative data of multiplier ideals. The chapter treats both algebraic and transcendental versions and proves the continuity and volume formulae needed in the big toric and non-Archimedean chapters.

Finally, \cref{chap:bdiv} translates the preceding singularity theory into the language of b-divisors. It associates with a positive current a nef b-divisor recording the regular parts of its pullbacks on all modifications, and it compares $\mathcal I$-equivalence with b-divisorial data. The mixed intersection theory of nef b-divisors developed there provides a geometric counterpart to the mixed volume theory of $\mathcal I$-good currents.

The chapters are arranged linearly. The comparison theory of \cref{chap:comp} is used throughout the part; \cref{chap:Igood,chap:trace} provide the main analytic tools; \cref{chap:testcurve,chap:Okounkov,chap:bdiv} develop three complementary languages for the same singularity theory: geodesic/test-curve language, Okounkov-body language and b-divisorial language. These languages are then used separately in the application chapters \cref{chap:toric_big,chap:NApotential,chap:Bergman}.
\end{partbacktext}
